
\section*{\underline{EXPOSÉ X}\\[3pt] \underline{\underline{FORMALISME DES INTERSECTIONS SUR LES SCHÉMAS ALGÉBRIQUES PROPRES}}}
\begin{center}
\underline{par O. Jussila}\\
(avec un appendice par A. Grothendieck)
\end{center}

\section*{1. Compatibilité des filtrations avec la loi de composition}
\[
K^{\bullet}(X)\times K_\bullet(X)\longrightarrow K_\bullet(X).
\]

\subsection*{1.1.}
Considérons d'abord un schéma noethérien $X$ et le groupe $K_\bullet(X)$ des classes de $\underline{O}_X$-modules cohérents. Pour tout $\underline{O}_X$-module cohérent $\underline{F}$ nous désignerons par $\cl_\bullet(\underline{F})$ sa classe dans $K_\bullet(X)$.

\textbf{Définition 1.1.1.} \uline{Pour tout entier} $n\in\mathbb Z$ \uline{on désigne par} $\Filt_n(X)$ \uline{le sous-groupe de} $K_\bullet(X)$ \uline{engendré par les classes} $\cl_\bullet(\underline{F})$ \uline{des} $\underline{O}_X$\uline{-modules cohérents} $\underline{F}$ \uline{tels que}
\[
\dim(\Supp(\underline{F}))\leq n.
\]
On définit ainsi une filtration croissante $(\Filt_n(X))_{n\in\mathbb Z}$ de $K_\bullet(X)$ telle que $\Filt_n(X)=0$ pour $n<0$.

\textbf{Proposition 1.1.2.} \uline{Soient} $X$ \uline{un schéma noethérien,} $\underline{F}$ \uline{un} $\underline{O}_X$\uline{-module cohérent tel que} $\dim(\Supp(\underline{F}))\leq n$\uline{, et} $X(n)$ ($n\in\mathbb Z$\uline{) le sous-ensemble} $\{x\in X\,;\,\dim(\overline{\{x\}})=n\}$ \uline{de} $X$\uline{. Pour tout} $x\in X$ \uline{on désigne par} $Z(x)$ \uline{le sous-schéma intègre fermé associé à} $\overline{\{x\}}$\uline{. Alors la longueur du} $\underline{O}_{X,x}$\uline{-module} $\underline{F}_x$ \uline{est finie pour tout} $x\in X(n)$ \uline{et non-nulle pour un nombre fini des} $x\in X(n)$ \uline{seulement, et on a}
